\appendix
\section{Phụ lục}

% ============================================================
% SLIDE S5
% ============================================================
\SectionSlide{A}{PHỤ LỤC}{Chi tiết toán học}{Dành cho người muốn đọc thêm}{%
  \SecItem{A1--A2}{Tucker và diễn giải}
  \SecItem{A3--A5}{Các bước cập nhật trong ALS}
  \SecItem{A6--A8}{MDL, Algorithm 2--4}}


% ============================================================
% SLIDE A1
% ============================================================
\begin{frame}{A1. Mode product và unfolding}{Phân rã Tucker}
\begin{columns}[T,onlytextwidth]
\begin{column}{0.50\textwidth}
\begin{mathcard}[Mode product]
  \small Với tensor $\mathcal G\in\mathbb R^{d_k\times d_l\times L}$ và ma trận $W\in\mathbb R^{d_k\times k}$:
  \[
  (\mathcal G\times_{\mathrm{key}}W)_{u\,j\,t}=\sum_{i=1}^{d_k} W_{iu}\,\mathcal G_{ijt}
  \]
  Kết quả cỡ $k\times d_l\times L$: chiều keyword ``phóng'' từ $d_k$ lên $k$.
\end{mathcard}
\vspace{2mm}
\begin{mathcard}[Áp dụng hai lần]
  \small $\widehat{\mathcal X}_d=\mathcal W^{(\mathrm{core})}\times_{\mathrm{key}}W^{(\mathrm{key})}\times_{\mathrm{loc}}W^{(\mathrm{loc})}$
  $\;\Leftrightarrow\;$
  $\hat x_{uvt}=\sum_{i}\sum_{j}W^{(\mathrm{key})}_{iu}W^{(\mathrm{loc})}_{jv}\,w_{ijt}$
\end{mathcard}
\end{column}
\begin{column}{0.47\textwidth}
\begin{mathcard}[Unfolding (matricization)]
  \small $\mathrm{unfold}(\mathcal X,\mathrm{key})\in\mathbb R^{k\times (l L)}$: xếp mọi ``sợi'' theo chiều keyword thành cột của một ma trận.
  Mode product trở thành nhân ma trận:
  \[
  \mathrm{unfold}(\mathcal G\times_{\mathrm{key}}W,\mathrm{key}) = W^{\top}\,\mathrm{unfold}(\mathcal G,\mathrm{key})
  \]
\end{mathcard}
\vspace{2mm}
\begin{intuition}[Kích thước, ví dụ Device]
  \small $\mathcal W^{(\mathrm{core})}$: $3\times2\times104$;\;
  $W^{(\mathrm{key})}$: $3\times8$;\;
  $W^{(\mathrm{loc})}$: $2\times50$\\
  $\Rightarrow\widehat{\mathcal X}_d$: $8\times50\times104$.
\end{intuition}
\end{column}
\end{columns}
\vfill
\papersource{Eq.~(2); Kolda \& Bader (2009), \emph{Tensor Decompositions and Applications}.}
\end{frame}

% ============================================================
% SLIDE A2
% ============================================================
\begin{frame}{A2. Diễn giải: chiếu hệ ẩn về từng chuỗi thật}{Interpretability of $\theta_d$}
\small Với từ khóa $u$ và địa điểm $v$ bất kỳ, kết hợp Eq.~(2) và (3):
\formula{\frac{d\,\mathcal X_{uv}}{dt}=
\underbrace{\sum_{i=1}^{d_k}\sum_{j=1}^{d_l}W^{(\mathrm{key})}_{iu}W^{(\mathrm{loc})}_{jv}}_{\text{chiếu đa tuyến tính}}
\underbrace{\Big(a_{ij}w_{ij}+\sum_{j'=1}^{d_l}d_{ijj'}(w_{ij'}-w_{ij})\Big)}_{\text{hệ phản ứng--khuếch tán ẩn}}}
\vspace{2mm}
\begin{columns}[T,onlytextwidth]
\begin{column}{0.48\textwidth}
\begin{intuition}[Ý nghĩa]
  Hệ ẩn mô tả $d_k\cdot d_l$ động thái quan trọng: $d_k$ nhóm từ khóa ẩn, $d_l$ nhóm vùng ẩn.
  Phép chiếu cho biết mỗi nhóm \textbf{đóng góp bao nhiêu} vào chuỗi thật của từ khóa $u$ tại địa điểm $v$.
\end{intuition}
\end{column}
\begin{column}{0.48\textwidth}
\begin{takeaway}[Vì sao vẫn diễn giải được]
  Hai ma trận chiếu \textbf{không âm}, nên không đảo chiều quan hệ giữa các động thái ẩn.
  ``Nhóm $j'$ lan sang nhóm $j$'' vẫn giữ nguyên chiều khi nhìn ở các nước thuộc hai nhóm đó.
\end{takeaway}
\end{column}
\end{columns}
\vfill
\papersource{§3.3, đoạn ``Interpretability of $\theta_d$'' sau Def.~2.}
\end{frame}

% ============================================================
% SLIDE A3
% ============================================================
\begin{frame}{A3. Cập nhật hệ phản ứng--khuếch tán $\theta$}{ModelEstimation --- §4.1.1}
\formula{\{\theta,w_0\}=\arg\min_{\theta',w_0'}\Big\|\mathcal X_c-\widehat{\mathcal X}_s-\widehat{\mathcal X}_o
-f_\theta(\theta',w_0',L_c)\times_{\mathrm{key}}W^{(\mathrm{key})}\times_{\mathrm{loc}}W^{(\mathrm{loc})}\Big\|}
\vspace{3mm}
\begin{itemize}\small
  \item Cố định $W^{(\mathrm{key})},W^{(\mathrm{loc})}$, $\theta_s$, $\widehat{\mathcal X}_o$; chỉ tối ưu $A$, $\mathcal D$ và trạng thái đầu $w_0$.
  \item $f_\theta$ là nghiệm số của ODE, phi tuyến theo tham số $\Rightarrow$ bài toán \textbf{bình phương tối thiểu phi tuyến}.
  \item Paper giải bằng \textbf{Levenberg--Marquardt}: lai giữa Gauss--Newton (gần nghiệm) và gradient descent (xa nghiệm).
  \item Số tham số: $|A|+|\mathcal D| = d_k d_l + d_k d_l^2$ $\Rightarrow$ chi phí $O(d_k d_l^2)$.
\end{itemize}
\vfill
\papersource{Eq.~(5), footnote 1; Algorithm 2 dòng 3--4; Appendix A.3.}
\end{frame}

% ============================================================
% SLIDE A4
% ============================================================
\begin{frame}{A4. Cập nhật $W^{(\mathrm{key})}$, $W^{(\mathrm{loc})}$ (multiplicative update)}{ModelEstimation --- §4.1.1}
\begin{columns}[T,onlytextwidth]
\begin{column}{0.56\textwidth}
\begin{mathcard}[Eq.~(6)]
  \small
  \[
  \begin{aligned}
  P^{(\mathrm{m})}&=\max\!\big(\epsilon,\;X_d^{(\mathrm{m})}\cdot G^{(\mathrm{m})\top}\big)\\
  Q^{(\mathrm{m})}&=\max\!\big(\epsilon,\;W^{(\mathrm{m})}\cdot G^{(\mathrm{m})}\cdot G^{(\mathrm{m})\top}\big)\\
  W^{(\mathrm{m})}&\leftarrow W^{(\mathrm{m})}\otimes P^{(\mathrm{m})}\oslash Q^{(\mathrm{m})}
  \end{aligned}
  \]
\end{mathcard}
\vspace{1mm}
\begin{mathcard}[Eq.~(7)]
  \small
  $X_d^{(\mathrm{m})}=\mathrm{unfold}(\mathcal X_c-\widehat{\mathcal X}_s-\widehat{\mathcal X}_o,\ \mathrm{m})$\\[1mm]
  $G^{(\mathrm{m})}=\mathrm{unfold}(\mathcal W^{(\mathrm{core})}\times_{\bar{\mathrm m}}W^{(\bar{\mathrm m})},\ \mathrm{m})$
\end{mathcard}
\end{column}
\begin{column}{0.41\textwidth}
\small
\begin{itemize}
  \item $\mathrm m\in\{\mathrm{key},\mathrm{loc}\}$; $\bar{\mathrm m}$ là mode còn lại.
  \item $\otimes$: nhân từng phần tử; $\oslash$: chia từng phần tử; $\epsilon$: hằng số nhỏ.
  \item Nhân với tỉ số của hai đại lượng không âm $\Rightarrow$ $W$ \textbf{luôn không âm}, không cần chiếu.
  \item Khởi tạo bằng \textbf{Nonnegative Tucker Decomposition} (NTD) trên $\widehat{\mathcal X}_d$.
  \item Chi phí $O(k d_k + l d_l)$.
\end{itemize}
\end{column}
\end{columns}
\vfill
\papersource{Eq.~(6)(7), dựa trên [24]; Appendix A.1, A.3.}
\end{frame}

% ============================================================
% SLIDE A5
% ============================================================
\begin{frame}{A5. Cập nhật seasonal factors (CP-ALS)}{ModelEstimation --- §4.1.2}
\formula{U^{(\mathrm m)}=\bigodot_{i\neq \mathrm m}S^{(i)},\qquad
V^{(\mathrm m)}=\mathop{\scalebox{1.3}{$\ast$}}_{i\neq \mathrm m}S^{(i)\top}S^{(i)},\qquad
S^{(\mathrm m)}=X_s^{(\mathrm m)}\,U^{(\mathrm m)}\,V^{(\mathrm m)\dagger}}
\vspace{3mm}
\begin{columns}[T,onlytextwidth]
\begin{column}{0.50\textwidth}
\small
\begin{itemize}
  \item $\mathrm m\in M=\{\mathrm{time},\mathrm{key},\mathrm{loc}\}$, cập nhật lần lượt từng mode.
  \item $\bigodot$: tích Khatri--Rao; $\ast$: tích Hadamard; $\dagger$: giả nghịch đảo Moore--Penrose.
  \item $X_s^{(\mathrm m)}=\mathrm{unfold}(\mathcal X_c-\widehat{\mathcal X}_d-\widehat{\mathcal X}_o,\ \mathrm m)$.
\end{itemize}
\end{column}
\begin{column}{0.46\textwidth}
\begin{intuition}[Ý tưởng]
  \small Cố định hai ma trận, ma trận thứ ba là nghiệm của một bài \textbf{bình phương tối thiểu tuyến tính}, có nghiệm đóng.
  Xoay vòng ba mode đến khi hội tụ.
\end{intuition}
\vspace{1mm}
{\small\color{DTMuted}Ghi chú: $\theta_s$ thực chất là phân rã CP (PARAFAC) rank $d_s$.\par}
\end{column}
\end{columns}
\vfill
\papersource{Eq.~(8), dựa trên [25] (Kolda \& Bader).}
\end{frame}

% ============================================================
% SLIDE A6
% ============================================================
\begin{frame}{A6. Hàm cost MDL đầy đủ}{Automatic tensor compression --- §4.2}
\begin{columns}[T,onlytextwidth]
\begin{column}{0.52\textwidth}
\begin{mathcard}[Model description cost]
  \scriptsize
  $\langle\mathcal F\rangle=\sum_{\Theta\in\mathcal F}\langle\Theta\rangle,\quad
  \langle\Theta\rangle=\langle\theta_d\rangle+\langle\theta_s\rangle+\langle\widehat{\mathcal X}_o\rangle$\\[1mm]
  $\langle\theta_d\rangle=\langle W^{(\mathrm{key})}\rangle+\langle W^{(\mathrm{loc})}\rangle+\langle A\rangle+\langle\mathcal D\rangle$\\[1mm]
  $\langle W^{(\mathrm{key})}\rangle=|W^{(\mathrm{key})}|(\log d_k+\log k+c_F)+\log^*|W^{(\mathrm{key})}|$\\[0.5mm]
  $\langle W^{(\mathrm{loc})}\rangle=|W^{(\mathrm{loc})}|(\log d_l+\log l+c_F)+\log^*|W^{(\mathrm{loc})}|$\\[0.5mm]
  $\langle A\rangle=|A|(\log d_k+\log d_l+c_F)+\log^*|A|$\\[0.5mm]
  $\langle\mathcal D\rangle=|\mathcal D|(\log d_k+2\log d_l+c_F)+\log^*|\mathcal D|$\\[0.5mm]
  $\langle S^{(\mathrm{time})}\rangle=|S^{(\mathrm{time})}|(\log d_s+\log L_c+c_F)+\log^*|S^{(\mathrm{time})}|$\\[0.5mm]
  {\color{DTMuted}(tương tự cho $S^{(\mathrm{key})}$, $S^{(\mathrm{loc})}$)}\\[0.5mm]
  $\langle\widehat{\mathcal X}_o\rangle=|\widehat{\mathcal X}_o|(\log n+\log k+\log l+c_F)+\log^*|\widehat{\mathcal X}_o|$
\end{mathcard}
\end{column}
\begin{column}{0.45\textwidth}
\begin{mathcard}[Data encoding cost]
  \small
  \[
  \langle\mathcal X\mid\mathcal F\rangle=\sum_{x\in\mathcal X}-\log_2 p_{\mu,\sigma}\big(x-\hat x_d-\hat x_s-\hat x_o\big)
  \]
\end{mathcard}
\vspace{1mm}
\small
\begin{itemize}
  \item $|\cdot|$: số phần tử khác 0.
  \item $c_F=32$ bit cho một số thực.
  \item $\log^*$: độ dài mã phổ quát cho số nguyên.
  \item Mỗi số khác 0 trả tiền cho \textbf{vị trí} (các $\log$) và \textbf{giá trị} ($c_F$).
\end{itemize}
\end{column}
\end{columns}
\vfill
\papersource{§4.2, Eq.~(9)(10), footnote 2.}
\end{frame}

% ============================================================
% SLIDE A7
% ============================================================
\begin{frame}[fragile]{A7. Algorithm 2 --- ModelEstimation}{Appendix A.2}
\begin{columns}[T,onlytextwidth]
\begin{column}{0.60\textwidth}
\begin{questioncard}
\small
\textbf{Input}: $\mathcal X_c$, $d_k,d_l,d_s$\quad \textbf{Output}: $\Theta$\par\vspace{1mm}
\begin{tabular}{@{}r@{\;\;}l@{}}
1: & $\{\theta_d,\theta_s\}\leftarrow$ Initialization$(\mathcal X_c,d_k,d_l,d_s)$;\; $\widehat{\mathcal X}_o\leftarrow0$\\
2: & \textbf{repeat}\\
3--4: & \quad $\theta\leftarrow\arg\min_{\theta^*}\|\mathcal X_c-\widehat{\mathcal X}_s-\widehat{\mathcal X}_o-f_{\theta_d}(\cdot)\|$ \hfill (Eq.~5)\\
5--6: & \quad $\mathcal W^{(\mathrm{core})}\leftarrow f_\theta(\theta,w_0,L)$\\
7: & \quad $\{W^{(\mathrm{key})},W^{(\mathrm{loc})}\}\leftarrow$ Update(\ldots) \hfill (Eq.~6)\\
8--9: & \quad $\{S^{(\mathrm{time})},S^{(\mathrm{key})},S^{(\mathrm{loc})}\}\leftarrow$ Update(\ldots) \hfill (Eq.~8)\\
10--11: & \quad $\widehat{\mathcal X}_o\leftarrow$ Sparsify$(\mathcal X_c-\widehat{\mathcal X}_d-\widehat{\mathcal X}_s)$\\
12: & \textbf{until} hội tụ\\
13: & \textbf{return} $\Theta=\{\theta_d,\theta_s,\widehat{\mathcal X}_o\}$
\end{tabular}
\end{questioncard}
\end{column}
\begin{column}{0.37\textwidth}
\begin{intuition}[Khởi tạo]
  \small STL tách $\mathcal X_c$ thành trend $\widehat{\mathcal X}_d$ và mùa vụ $\widehat{\mathcal X}_s$.
  PARAFAC trên $\widehat{\mathcal X}_s$ $\to$ $S^{(\cdot)}$; NTD trên $\widehat{\mathcal X}_d$ $\to$ $W^{(\mathrm{key})},W^{(\mathrm{loc})}$.
\end{intuition}
\vspace{2mm}
\begin{takeaway}[Lemma 1]
  \small Mỗi vòng $O(d_k d_l^2+k d_k+l d_l)$; số vòng là hằng nhỏ
  $\Rightarrow$ tổng $O(d_k d_l^2+k d_k+l d_l)$.\\
  Không có $L$ hay $t_c$ trong độ phức tạp.
\end{takeaway}
\end{column}
\end{columns}
\vfill
\papersource{Appendix A.1--A.3, Algorithm 2, Lemma 1.}
\end{frame}

% ============================================================
% SLIDE A8
% ============================================================
\begin{frame}{A8. ModelUpdate và RankUpdate}{Appendix A.4}
\begin{columns}[T,onlytextwidth]
\begin{column}{0.49\textwidth}
\begin{mathcard}[ModelUpdate (Algorithm 3)]
  \small
  \begin{enumerate}
    \item Tính cost khi mô hình hóa $\mathcal X_c$ bằng các $\Theta\in\mathcal F$ hiện có.
    \item Tính cost khi thêm ứng viên $\Theta'$: $\langle\mathcal X;\mathcal F\cup\{\Theta'\}\rangle$.
    \item Nếu rẻ hơn: $\mathcal F\leftarrow\mathcal F\cup\{\Theta'\}$, dùng $\Theta'$ cho giai đoạn hiện tại.
    \item Gọi RankUpdate để điều chỉnh $d_k,d_l,d_s$.
  \end{enumerate}
\end{mathcard}
\end{column}
\begin{column}{0.48\textwidth}
\begin{mathcard}[RankUpdate (Algorithm 4)]
  \small
  \begin{itemize}
    \item Ứng viên: $(d_k\pm1,d_l,d_s)$, $(d_k,d_l\pm1,d_s)$, $(d_k,d_l,d_s\pm1)$.
    \item Với mỗi ứng viên: chạy ModelEstimation, tính cost MDL.
    \item Giữ rank cho cost nhỏ nhất (kể cả rank cũ).
  \end{itemize}
\end{mathcard}
\vspace{2mm}
\begin{intuition}[Vì sao chỉ đổi $\pm1$?]
  \small Giả định cấu trúc tối ưu thay đổi \textbf{từ từ} trong luồng $\Rightarrow$ không phải grid search lại toàn bộ, giữ chi phí thấp. Đây là nguồn của các ``gai'' trong Figure 4.
\end{intuition}
\end{column}
\end{columns}
\vfill
\papersource{§4.3, Appendix A.4, Algorithm 3--4.}
\end{frame}
